\documentclass[10pt, a4paper]{article}

% --- Packages ---
\usepackage[top=2cm, bottom=2cm, left=2cm, right=2cm]{geometry}
\usepackage{graphicx} % For images
\usepackage{amsmath}  % For math/equations
\usepackage{hyperref} % For clickable links
\usepackage{caption}  % For styling captions
\usepackage{titlesec} % For controlling section title sizes
\usepackage{float}
\usepackage{xcolor} % Required for defining and using colors
\usepackage{hyperref} % Keep this as your last \usepackage
\usepackage{color}
\usepackage{soul}

% ברירת המחדל של המרקור היא צהוב. אם תרצי לשנות למשל לטורקיז/תכלת, הוסיפי גם את השורה הזו:
% \sethlcolor{cyan}

\hypersetup{
    colorlinks=true, % Removes the default boxes and colors the text instead
    linkcolor=blue,  % Color for internal links (e.g., Figure and Table references)
    citecolor=blue,  % Color for bibliographical citations
    urlcolor=blue    % Color for external web URLs
}
% --- Document Formatting ---
\setlength{\parindent}{0pt}
\setlength{\parskip}{0.5em}
\captionsetup{font=scriptsize, labelfont=bf}

% --- Customize Title Sizes ---
\titleformat{\section}
  {\large\bfseries}{\thesection}{1em}{}
\titleformat{\subsection}
  {\normalsize\bfseries}{\thesubsection}{1em}{}

\usepackage{fancyhdr}
\pagestyle{fancy}
\fancyhf{} % מנקה את הגדרות ברירת המחדל
\fancyhead[L]{\footnotesize Weekly Summary \#3 \today} % יופיע למעלה בשמאל
\fancyhead[R]{\footnotesize Rachel Broner $\vert$ Miri Adler's Lab} % יופיע למעלה בימין
\fancyfoot[C]{\thepage} % מספר עמוד למטה באמצע

% \usepackage{mathptmx}

\usepackage{setspace}
\onehalfspacing % For 1.5 line spacing

% --- One place to change how a figure is drawn ---
\newcommand{\onefig}[2]{%
  \begin{figure}[H]\centering
    \includegraphics[width=0.72\textwidth]{summary_downloads/#1.pdf}
    \caption{#2}
  \end{figure}}
\newcommand{\widefig}[2]{%
  \begin{figure}[H]\centering
    \includegraphics[width=\textwidth]{summary_downloads/#1.pdf}
    \caption{#2}
  \end{figure}}
% Tall figures must be bounded on BOTH sides: with only a height, a wide figure is
% scaled until it runs off the page.
\newcommand{\tallfig}[2]{%
  \begin{figure}[H]\centering
    \includegraphics[width=\textwidth, height=0.80\textheight, keepaspectratio]{summary_downloads/#1.pdf}
    \caption{#2}
  \end{figure}}
\newcommand{\twofig}[3]{%
  \begin{figure}[H]\centering
    \includegraphics[width=0.49\textwidth]{summary_downloads/#1.pdf}\hfill
    \includegraphics[width=0.49\textwidth]{summary_downloads/#2.pdf}
    \caption{#3}
  \end{figure}}

\begin{document}

\begin{center}
    {\Large \textbf{Weekly Summary \#6 $\vert$ Association Rule Mining}} \\[0.5em]
    {\large Part 2 $\vert$ FOVs as vectors in rule space} \\[0.5em]
    \textit{GVHD Analysis} $\vert$ Week of: \today
\end{center}
\vspace{0.2cm}
\hrule
\vspace{0.5cm}

% The one line here that is about this dataset rather than about the method.
The data are MIBI images of gut biopsies taken after stem-cell transplant and of healthy
controls, one field of view (FOV) per image, with every cell typed and placed; the second
control group (\texttt{Control\_S}) is left out of everything below.

Each FOV is turned into a vector: one entry per rule, holding that rule's lift in that
FOV, and the neutral value where the rule was not found. PCA then places the FOVs so that
FOVs with similar rules sit near each other. Nothing about the group an FOV belongs to
enters the calculation -- the colours are added afterwards, to see whether the groups came
apart on their own.

The three sections below are built the same way: the PCA on the positive rules alone, then
the same PCA on positive and negative rules together with what drives it, and then the
checks asking whether a separation is about where the cells sit or only about which cells
are there.

% ===========================================================================
\section{All FOVs: colon against duodenum}
% ===========================================================================

Every FOV of both organs in one PCA, coloured by organ.

\subsection{Positive rules only}

\textit{Do the two organs come apart when only attracting pairs are used?}

\onefig{pos_all_pca_organ}{%
One dot per FOV, placed by its rules and coloured by organ. Only rules with lift above 1
take part, so every column says how much more often two cell types are neighbours than
they would be by chance. The axis labels give the share of the total variation each
component carries.}

\subsection{Positive and negative rules}

\textit{Does adding the avoiding pairs change where the FOVs land?}

\onefig{posneg_all_pca_organ}{%
The same FOVs with rules of both signs: cell types that attract, and cell types that avoid
each other. Circled FOVs are the ones drawn as tissue below; the blue square marks the
closest pair of FOVs belonging to different groups.}

\textit{Two FOVs that land almost on top of each other -- do they look alike as tissue?}

\widefig{posneg_all_pair_fovs}{%
The boxed pair from the figure above, drawn as maps, one dot per cell coloured by cell
type. These two are as close as any two FOVs from different groups get.}

\textit{Walking along the first component, does the tissue change?}

\widefig{posneg_all_fovs_pc1}{%
Five FOVs spread across the range of the first component, from one end to the other. If
the component describes something real about the tissue, it should show as a gradient
across these maps.}

\textit{Which rules pull each component?}

\twofig{posneg_all_loadings_pc1}{posneg_all_loadings_pc2}{%
The rules with the largest weight on the first component (left) and on the second (right).
A bar to the right means the rule pushes an FOV towards the positive end of that
component, to the left means the negative end.}

\subsection{Is it the rules, or how common each cell type is?}

\textit{Does a component simply follow how much of some cell type an FOV holds?}

\onefig{abund_all_abundance_correlation}{%
For every cell type, how closely its share of the cells in an FOV moves together with each
component. The measure is Spearman rho, a correlation of ranks, so a single extreme FOV
cannot drive it. Only the strongest cell types are drawn; inside the grey band a value is
too small to mean anything at this number of FOVs.}

\textit{Seen on the PCA itself, does that cell type line up with a direction?}

\twofig{abund_all_abundance_top1}{abund_all_abundance_top2}{%
The same PCA coloured by the share of the two cell types whose correlation with the
components was strongest. A smooth colour gradient across the plot means the component is
largely tracking how common that cell type is.}

\subsection{Rules against cell counts}

\textit{Would counting the cells have placed the FOVs in the same way?}

\tallfig{abund_all_panel}{%
Top row: the FOVs placed by their rules (left), and by their cell composition alone
(right), with no rules involved. Each row below repeats the pair with one
composition-linked cell type removed -- its rules dropped on the left, its column dropped
on the right. Removing them one at a time, rather than all together, shows whether any
single cell type is carrying the separation.}

\subsection{Splitting the FOVs by what they are made of}

\textit{If the FOVs are sorted into two groups purely by composition, do those groups turn
out to be the groups we care about?}

\widefig{abund_all_groups}{%
The FOVs split in two by how much they hold of the cell types tied to this PCA (k-means on
the square root of each share, so that neither the most common nor the rarest cell type
dominates). Left: how the two groups fall against the labels -- a group holding only one
label means the cell counts alone already separate those FOVs. Right: what each group is
made of on average.}

\textit{What does each individual FOV in those groups look like?}

\widefig{abund_all_group_bars}{%
One bar per FOV, sorted into the two groups. The coloured segments are the cell types the
split was made on, as a share of all the cells in that FOV; the grey block on top is
everything else, so the height of the colour shows how much of the tissue the split
actually used. The strips above give each FOV's group and its label.}

\textit{Beyond those few cell types, do the two groups differ in everything else as well?}

\tallfig{abund_all_group_heatmap}{%
Every cell type against every FOV. Each cell type is standardised: red means an FOV holds
more of it than a typical FOV here, blue means less, so a row can be read along but not
compared with another row. Rows are ordered by the gap between the two groups, drawn again
as bars on the right; starred cell types are the ones the split was made on.}

\textit{Is the separation about which rules were found, or about how strong they are?}

\widefig{abund_all_presence}{%
Three matrices over the same FOVs and the same rules. Left: the lift values. Middle: the
same rules reduced to found or not found, the values thrown away. Right: whether the FOV
held enough of both of a rule's cell types for that rule to be findable at all -- a matrix
carrying no information about where any cell sits. If the third separates the groups as
well as the other two, the separation does not need the arrangement.}

\clearpage
% ===========================================================================
\section{Colon}
% ===========================================================================

\subsection{Control against Severe}

\onefig{pos_cs_colon_pca_pathological_score}{%
Positive rules only. One dot per FOV, coloured by pathological score.}

\onefig{posneg_cs_colon_pca_pathological_score}{%
Positive and negative rules. Circled FOVs are drawn as tissue below; the blue square marks
the closest pair from different groups.}

\widefig{posneg_cs_colon_pair_fovs}{%
The boxed pair, drawn as maps: one dot per cell, coloured by cell type.}

\widefig{posneg_cs_colon_fovs_pc1}{%
Five FOVs spread across the range of the first component.}

\twofig{posneg_cs_colon_loadings_pc1}{posneg_cs_colon_loadings_pc2}{%
The rules with the largest weight on the first component (left) and on the second (right).}

\onefig{abund_cs_colon_abundance_correlation}{%
How closely each cell type's share of the cells follows each component. Inside the grey
band a value is too small to mean anything at this number of FOVs.}

\twofig{abund_cs_colon_abundance_top1}{abund_cs_colon_abundance_top2}{%
The same PCA coloured by the share of the two cell types most tied to the components.}

\tallfig{abund_cs_colon_panel}{%
Top row: the FOVs placed by their rules, and by their cell composition alone. Each row
below repeats it with one composition-linked cell type left out -- its rules on the left,
its column on the right.}

\widefig{abund_cs_colon_groups}{%
The FOVs split in two by their composition of the cell types tied to this PCA. Left: how
those groups fall against the labels. Right: what each group is made of.}

\widefig{abund_cs_colon_group_bars}{%
One bar per FOV, sorted into the two groups: the share of each cell type the split was made
on, with everything else in grey.}

\tallfig{abund_cs_colon_group_heatmap}{%
Every cell type against every FOV, each cell type standardised. Rows are ordered by the gap
between the two groups, drawn again as bars on the right.}

\widefig{abund_cs_colon_presence}{%
The lift values, the same rules as found or not found, and whether each rule could have
been found at all from the cell counts alone.}

\clearpage
\subsection{Mild against Severe}

\onefig{pos_ms_colon_pca_pathological_score}{%
Positive rules only. One dot per FOV, coloured by pathological score.}

\onefig{posneg_ms_colon_pca_pathological_score}{%
Positive and negative rules. Circled FOVs are drawn as tissue below; the blue square marks
the closest pair from different groups.}

\widefig{posneg_ms_colon_pair_fovs}{%
The boxed pair, drawn as maps: one dot per cell, coloured by cell type.}

\widefig{posneg_ms_colon_fovs_pc1}{%
Five FOVs spread across the range of the first component.}

\clearpage
% ===========================================================================
\section{Duodenum}
% ===========================================================================

\subsection{Control against Severe}

\onefig{pos_cs_duodenum_pca_pathological_score}{%
Positive rules only. One dot per FOV, coloured by pathological score.}

\onefig{posneg_cs_duodenum_pca_pathological_score}{%
Positive and negative rules. Circled FOVs are drawn as tissue below; the blue square marks
the closest pair from different groups.}

\widefig{posneg_cs_duodenum_pair_fovs}{%
The boxed pair, drawn as maps: one dot per cell, coloured by cell type.}

\widefig{posneg_cs_duodenum_fovs_pc1}{%
Five FOVs spread across the range of the first component.}

\twofig{posneg_cs_duodenum_loadings_pc1}{posneg_cs_duodenum_loadings_pc2}{%
The rules with the largest weight on the first component (left) and on the second (right).}

\onefig{abund_cs_duodenum_abundance_correlation}{%
How closely each cell type's share of the cells follows each component. Inside the grey
band a value is too small to mean anything at this number of FOVs.}

\twofig{abund_cs_duodenum_abundance_top1}{abund_cs_duodenum_abundance_top2}{%
The same PCA coloured by the share of the two cell types most tied to the components.}

\tallfig{abund_cs_duodenum_panel}{%
Top row: the FOVs placed by their rules, and by their cell composition alone. Each row
below repeats it with one composition-linked cell type left out -- its rules on the left,
its column on the right.}

\widefig{abund_cs_duodenum_groups}{%
The FOVs split in two by their composition of the cell types tied to this PCA. Left: how
those groups fall against the labels. Right: what each group is made of.}

\widefig{abund_cs_duodenum_group_bars}{%
One bar per FOV, sorted into the two groups: the share of each cell type the split was made
on, with everything else in grey.}

\tallfig{abund_cs_duodenum_group_heatmap}{%
Every cell type against every FOV, each cell type standardised. Rows are ordered by the gap
between the two groups, drawn again as bars on the right.}

\widefig{abund_cs_duodenum_presence}{%
The lift values, the same rules as found or not found, and whether each rule could have
been found at all from the cell counts alone.}

\clearpage
\subsection{Mild against Severe}

\onefig{pos_ms_duodenum_pca_pathological_score}{%
Positive rules only. One dot per FOV, coloured by pathological score.}

\onefig{posneg_ms_duodenum_pca_pathological_score}{%
Positive and negative rules. Circled FOVs are drawn as tissue below; the blue square marks
the closest pair from different groups.}

\widefig{posneg_ms_duodenum_pair_fovs}{%
The boxed pair, drawn as maps: one dot per cell, coloured by cell type.}

\widefig{posneg_ms_duodenum_fovs_pc1}{%
Five FOVs spread across the range of the first component.}

\clearpage
% ===========================================================================
\section{Do a patient's FOVs sit together?}
% ===========================================================================

\textit{FOVs from one patient are not independent -- how much of the spread is
just which patient the tissue came from?}

\widefig{posneg_all_patient_spread}{%
One dot per patient, over all FOVs of both organs. Across is how far that patient's
own FOVs sit from each other, divided by how far any two FOVs sit from each other, so
1.0 means the patient's FOVs are no more alike than any others and left of it means
they cluster. Up is how many FOVs the patient has; the dots are nudged off the whole
numbers so that patients sharing a count do not stack into one line.}

\end{document}
